%! Equivalent Representations of Polynomial Expressions — Unit 2, Lesson 2.6 — Guided Notes
\documentclass[10pt]{article}
\usepackage{precalculus-article}
\usepackage{precalculus-boxes}
\newcommand{\TallMath}[1]{$\displaystyle #1\rule[-1.4em]{0pt}{3.2em}$}

\begin{document}

\pageheader{Unit 2, Lesson 2.6}{Guided Notes}

\vspace{0.12in}

\begin{objectivebox}
By the end of this lesson, I will be able to\ldots
\begin{enumerate}[itemsep=2pt]
  \item Decide which form of a polynomial --- \blank{2cm} or
        \blank{2cm} --- answers a given question.
  \item Rewrite a polynomial in factored form as a polynomial in standard form
        by \blank{2.8cm} it out.
  \item Use one row of \blank{2.6cm} Triangle to expand $(x+c)^n$.
  \item Divide one polynomial by another and write
        $f(x) = g(x)q(x) + r(x)$, where the \blank{2.4cm} has lower degree
        than the divisor.
\end{enumerate}
\end{objectivebox}

\vspace{0.1in}

\boxguard
\begin{vocabbox}
Fill in each term as we define it together.
\par\vspace{2pt}
\noindent\textbf{\textcolor{plum}{Standard form:}}\\[1pt]
\writeline\\[3pt]
\noindent\textbf{\textcolor{plum}{Factored form:}}\\[1pt]
\writeline\\[3pt]
\noindent\textbf{\textcolor{plum}{Pascal's Triangle:}}\\[1pt]
\writeline\\[3pt]
\noindent\textbf{\textcolor{plum}{Binomial theorem:}}\\[1pt]
\writeline\\[3pt]
\noindent\textbf{\textcolor{plum}{Polynomial long division:}}\\[1pt]
\writeline\\[3pt]
\noindent\textbf{\textcolor{plum}{Quotient and remainder:}}\\[1pt]
\writeline\\[3pt]
\noindent\textbf{\textcolor{plum}{Division algorithm:}}\\[1pt]
\writeline
\end{vocabbox}

\vspace{0.1in}

\boxguard[34]
\begin{hookbox}
These two lines describe the \textit{same} function:
\[
  p(x) = (x+1)(x-2)(x-3) \qquad\text{and}\qquad p(x) = x^3 - 4x^2 + x + 6 .
\]

\begin{center}
\begin{tikzpicture}[x=0.9cm, y=0.22cm]
  \draw[linegray, very thin, xstep=1, ystep=2] (-1.6,-8) grid (3.6,6);
  \draw[->, thick] (-1.9,0) -- (3.9,0) node[below right] {\small $x$};
  \draw[->, thick] (0,-8.6) -- (0,6.8) node[above] {\small $p(x)$};
  \foreach \x in {-1,1,2,3} \node[below] at (\x,-0.4) {\small \x};
  \foreach \y in {-8,-6,-4,-2,2,4,6} \node[left] at (0,\y) {\small \y};
  \draw[very thick, plum] plot [smooth] coordinates
    {(-1.5,-7.88) (-1.3,-4.26) (-1.1,-1.27) (-1,0) (-0.8,2.13)
     (-0.5,4.38) (0,6) (0.5,5.63) (1,4) (1.5,1.88) (2,0)
     (2.5,-0.88) (2.8,-0.61) (3,0) (3.3,1.68) (3.5,3.38)};
  \fill[plum] (-1,0) circle (2.5pt);
  \fill[plum] (2,0) circle (2.5pt);
  \fill[plum] (3,0) circle (2.5pt);
\end{tikzpicture}
\end{center}

Two questions: \textit{Where does the graph cross the $x$-axis?} and
\textit{What is $\lim_{x \to -\infty} p(x)$?} Which line would you use for
each --- and why is it not the same line twice?

\writelines{2}
\end{hookbox}

\vspace{0.1in}

\boxguard[30]
\begin{notesbox}{1. Which Form Answers Which Question?}
Neither form is ``more simplified'' than the other. They are the same
function, and each one puts different information on the surface.

\smallskip
\renewcommand{\arraystretch}{1.5}
\begin{center}
\begin{tabular}{|p{0.30\linewidth}|C{2.4cm}|p{0.36\linewidth}|}
\hline
\textbf{The question} & \textbf{Which form?} & \textbf{What that form shows} \\
\hline
Where does the graph cross the $x$-axis? & \blank{1.8cm} &
Each linear factor $(x-a)$ hands over the real zero $a$. \\
\hline
What is $\lim_{x \to \pm\infty} p(x)$? & \blank{1.8cm} &
The leading term $a_nx^n$ decides both ends. \\
\hline
What is the $y$-intercept? & \blank{1.8cm} &
The constant term $a_0$ \textit{is} $p(0)$. \\
\hline
What is the degree? & \blank{1.8cm} &
Add the exponents on the factors, or read the highest power. \\
\hline
\end{tabular}
\end{center}
\smallskip

\textbf{The hook, answered.} Use $p(x) = (x+1)(x-2)(x-3) = x^3-4x^2+x+6$.

\begin{itemize}[itemsep=4pt]
  \item $x$-intercepts: \blank{3.6cm} \quad (from the \blank{2cm} form)
  \item $\lim_{x \to -\infty} p(x) = $ \blank{1.6cm} \quad (from the
        \blank{2cm} form --- leading term $x^3$)
  \item $y$-intercept: \blank{1.6cm} \quad Check it the other way:
        $p(0) = (1)(-2)(-3) = $ \blank{1.2cm}
\end{itemize}

\textbf{The habit to build:} read the question first, \textit{then} choose the
form. Changing forms is work --- do it on purpose.
\end{notesbox}

\vspace{0.1in}

\boxguard[30]
\begin{notesbox}{2. Factored $\rightarrow$ Standard, and the Binomial Theorem}
Going from factored to standard form is just distributing, two factors at a
time.

\textbf{Example 1.} Expand $p(x) = (x+1)(x-2)(x-3)$.

\begin{work}
  p(x) &= (x+1)(x-2)(x-3) \\
       &= \big(x^2 - x - 2\big)(x-3) \\
       &= x^3 - 3x^2 - x^2 + 3x - 2x + 6 \\
       &= x^3 - 4x^2 + x + 6
\end{work}

\medskip
\textbf{The shortcut for $(x+c)^n$.} Multiplying four copies of $(x+2)$ by
hand is slow. \textbf{Pascal's Triangle} gives the coefficients directly ---
each entry is the sum of the two directly above it.

\smallskip
\begin{center}
\renewcommand{\arraystretch}{1.2}
\begin{tabular}{r c}
row $0$: & $1$ \\
row $1$: & $1 \quad 1$ \\
row $2$: & $1 \quad 2 \quad 1$ \\
row $3$: & $1 \quad 3 \quad 3 \quad 1$ \\
row $4$: & $1 \quad 4 \quad 6 \quad 4 \quad 1$ \\
row $5$: & $1 \quad 5 \quad 10 \quad 10 \quad 5 \quad 1$ \\
\end{tabular}
\end{center}
\smallskip

\textbf{Binomial theorem.} To expand $(a+b)^n$, use row \blank{1cm} of the
triangle. Across the terms, the power of $a$ \blank{1.6cm} from $n$ down to
$0$ while the power of $b$ \blank{1.6cm} from $0$ up to $n$.

\textbf{Example 2.} Expand $(x+2)^4$. Row 4 is $1, 4, 6, 4, 1$.

\begin{work}
  (x+2)^4 &= x^4 + 4x^3(2) + 6x^2(2)^2 + 4x(2)^3 + (2)^4 \\
          &= x^4 + 8x^3 + 24x^2 + 32x + 16
\end{work}

Sanity check against Lesson 2.3: $(x+2)^4$ has exactly one real zero,
$x = $ \blank{1.4cm}, with multiplicity \blank{1.2cm} --- and the expanded
form has leading term \blank{1.4cm} and constant term \blank{1.4cm}.
\end{notesbox}

\vspace{0.1in}

\boxguard[30]
\begin{notesbox}{3. Long Division and the Division Algorithm}
Numbers first: $17 \div 5$ is $3$ with remainder $2$, so
$17 = 5 \cdot 3 + 2$ --- and the remainder is \textit{smaller} than the
divisor. Polynomials work the same way, with ``smaller'' meaning
\textbf{lower degree}.

\medskip
\textbf{The division algorithm.} If $f$ is divided by $g$, then
\[
  f(x) = g(x)\,q(x) + r(x), \qquad \deg r < \deg g,
\]
where $q$ is the \blank{2.2cm} and $r$ is the \blank{2.4cm}.

\medskip
\textbf{The routine:} divide $\rightarrow$ multiply $\rightarrow$ subtract
$\rightarrow$ bring down. Repeat until what is left has lower degree than the
divisor.

\textbf{Example 3.} Divide $2x^3 - 3x^2 - 5x + 6$ by $x - 2$.

\begin{work}
  &\begin{array}{r@{\;}rrrr}
              &  2x^2 &   +x  &  -3  &     \\ \cline{2-5}
   x-2\,\big) &  2x^3 & -3x^2 & -5x  & +6  \\
   -          &  2x^3 & -4x^2 &      &     \\ \cline{2-5}
              &       &   x^2 & -5x  & +6  \\
   -          &       &   x^2 & -2x  &     \\ \cline{2-5}
              &       &       & -3x  & +6  \\
   -          &       &       & -3x  & +6  \\ \cline{2-5}
              &       &       &      &  0
  \end{array}
\end{work}

Quotient: \blank{3.6cm} \qquad Remainder: \blank{1.2cm}

\medskip
\textbf{Cash it in against Lesson 2.3.} The remainder is $0$, so
$2x^3 - 3x^2 - 5x + 6 = (x-2)(2x^2 + x - 3)$ --- which says $(x-2)$ is a
\blank{1.8cm} of $f$, which says $x = 2$ is a \blank{1.6cm} of $f$. All three
statements are the same statement:
\[
  r = 0 \iff (x-a) \text{ is a factor of } f \iff f(a) = 0 .
\]
So long division is how you \textbf{test} a proposed factor, and how you
\textbf{reduce} a polynomial once one zero is known.

\medskip
\textbf{Example 4.} Divide $x^3 - 4x^2 + 2x + 5$ by $x - 3$. This one does
\textit{not} come out even.

\begin{work}
  &\begin{array}{r@{\;}rrrr}
              &   x^2 &   -x  &  -1  &     \\ \cline{2-5}
   x-3\,\big) &   x^3 & -4x^2 & +2x  & +5  \\
   -          &   x^3 & -3x^2 &      &     \\ \cline{2-5}
              &       &  -x^2 & +2x  & +5  \\
   -          &       &  -x^2 & +3x  &     \\ \cline{2-5}
              &       &       &  -x  & +5  \\
   -          &       &       &  -x  & +3  \\ \cline{2-5}
              &       &       &      &  2
  \end{array}
\end{work}

Write the result in division-algorithm form:

\smallskip
$f(x) = (x-3)\,\cdot$ \blank{3.8cm} $+$ \blank{1.2cm}

\smallskip
Degree check: $\deg r = $ \blank{1cm} and $\deg g = $ \blank{1cm}, so
$\deg r < \deg g$. \quad Is $x-3$ a factor of $f$? \blank{1.6cm}
\end{notesbox}

\vspace{0.1in}

\boxguard[30]
\begin{notesbox}{4. Both Forms, One Context}
An open-top box is made from a $12$ in.\ by $16$ in.\ sheet of cardboard by
cutting a square of side $x$ from each corner and folding up the sides. The
volume is
\[
  V(x) = x(12-2x)(16-2x).
\]

\textbf{Expand it into standard form.}

\begin{work}
  V(x) &= x(12-2x)(16-2x) \\
       &= x\big(192 - 24x - 32x + 4x^2\big) \\
       &= x\big(4x^2 - 56x + 192\big) \\
       &= 4x^3 - 56x^2 + 192x
\end{work}

Now let each form do its own job.

\begin{itemize}[itemsep=4pt]
  \item \textbf{Factored form.} $V(x) = 0$ when $x = $ \blank{3.6cm}.
  \item A real box needs $x > 0$ \textit{and} both side lengths positive, so
        the domain in context is $0 < x < $ \blank{1.2cm}. The zero
        $x = $ \blank{1cm} is algebraically correct and physically
        impossible.
  \item \textbf{Standard form.} Leading term: \blank{1.8cm}, so
        $\lim_{x \to \infty} V(x) = $ \blank{1.4cm}.
  \item In context that limit is \blank{3.6cm} --- the model only describes
        the box for $0 < x < 6$.
\end{itemize}
\end{notesbox}

\vspace{0.1in}

\boxguard
\begin{practicebox}
\begin{enumerate}[itemsep=8pt]
  \item Let $p(x) = (x+4)(x-1)(x-2)$. Write $p$ in standard form, then read
        off the $y$-intercept and the end behavior.

        \begin{work}
          p(x) &= (x+4)(x-1)(x-2) \\
               &= \big(x^2 + 3x - 4\big)(x-2) \\
               &= x^3 - 2x^2 + 3x^2 - 6x - 4x + 8 \\
               &= x^3 + x^2 - 10x + 8
        \end{work}

        Standard form: \blank{6cm}

        $y$-intercept: \blank{1.4cm} \qquad
        $\lim_{x \to -\infty} p(x) = $ \blank{1.6cm}

  \item Divide $x^3 - 5x + 3$ by $x - 2$. Watch the missing $x^2$ term ---
        write it as $0x^2$.

        \begin{work}
          &\begin{array}{r@{\;}rrrr}
                      &   x^2 &  +2x  &  -1  &     \\ \cline{2-5}
           x-2\,\big) &   x^3 & +0x^2 & -5x  & +3  \\
           -          &   x^3 & -2x^2 &      &     \\ \cline{2-5}
                      &       &  2x^2 & -5x  & +3  \\
           -          &       &  2x^2 & -4x  &     \\ \cline{2-5}
                      &       &       &  -x  & +3  \\
           -          &       &       &  -x  & +2  \\ \cline{2-5}
                      &       &       &      &  1
          \end{array}
        \end{work}

        Quotient: \blank{3.6cm} \qquad Remainder: \blank{1.2cm}

        Is $x-2$ a factor of $x^3 - 5x + 3$? \quad \blank{1.6cm}
\end{enumerate}
\end{practicebox}

\end{document}
